% GENERATED FILE. Edit canonical lessons, then run npm run generate:books.
% canonical-source-hash: fnv1a64:ef712c287f5bc919
% canonical-lessons: TA-C214-mor, TA-C214-palaccaru, TA-C214-ilanir, TA-C214-mavu, TA-C214-rotti

\chapter{Buttermilk, Fruit juice, Tender coconut water, Flour, Bread}
\label{ch:ta-buttermilk-fruit-juice-tender-coconut-water-flour-bread}
\hlchaptermodality{214}

\begin{chapteropening}
\textbf{By the end of this chapter:} \emph{I can say buttermilk, fruit juice, tender coconut water, flour and bread.}

Complete the last lesson of chapter 214: I can say buttermilk, fruit juice, tender coconut water, flour and bread.
\end{chapteropening}

\section[\texorpdfstring{\ta{மோர்} (mōr)}{மோர் (mōr)}]{\ta{மோர்} (mōr) — buttermilk}

\label{lesson:TA-C214-mor}



\begin{quote}
Before the new one: say the Tamil for toothpaste, then the Tamil for a toothbrush.
\end{quote}

\subsection*{You'll want to know: \ta{மோர்}}
\textbf{\ta{மோர்}} — \emph{mōr} — \textquotedblleft{}buttermilk\textquotedblright{}.

\begin{grammarlens}[title={What you've built}]
One more word for this chapter.
\end{grammarlens}

\subsection*{Your turn}
\begin{itemize}
\raggedright
  \item \emph{Say it:} \emph{mōr}
  \item \emph{Say it:} \emph{mōr}, once more
  \item \emph{Say it:} say \emph{mōr}
  \item \emph{Recall:} say the Tamil for toothpaste, then the Tamil for a toothbrush, then say \emph{mōr} again
\end{itemize}

\begin{tcolorbox}[breakable,colback=teal!4,colframe=teal!35!black,title={Before you move on}]
What does \ta{மோர்} mean? (\textquotedblleft{}Buttermilk\textquotedblright{}.) Say it once more.
\end{tcolorbox}
\section[\texorpdfstring{\ta{பழச்சாறு} (paḻaccāṟu)}{பழச்சாறு (paḻaccāṟu)}]{\ta{பழச்சாறு} (paḻaccāṟu) — fruit juice}

\label{lesson:TA-C214-palaccaru}



\begin{quote}
Before the new one: say the Tamil for a toothbrush, then the Tamil for buttermilk.
\end{quote}

\subsection*{You'll want to know: \ta{பழச்சாறு}}
\textbf{\ta{பழச்சாறு}} — \emph{paḻaccāṟu} — \textquotedblleft{}fruit juice\textquotedblright{}.

\begin{grammarlens}[title={What you've built}]
One more word for this chapter.
\end{grammarlens}

\subsection*{Your turn}
\begin{itemize}
\raggedright
  \item \emph{Say it:} \emph{paḻaccāṟu}
  \item \emph{Say it:} \emph{paḻaccāṟu}, once more
  \item \emph{Say it:} say \emph{paḻaccāṟu}
  \item \emph{Recall:} say the Tamil for a toothbrush, then the Tamil for buttermilk, then say \emph{paḻaccāṟu} again
\end{itemize}

\begin{tcolorbox}[breakable,colback=teal!4,colframe=teal!35!black,title={Before you move on}]
What does \ta{பழச்சாறு} mean? (\textquotedblleft{}Fruit juice\textquotedblright{}.) Say it once more.
\end{tcolorbox}
\section[\texorpdfstring{\ta{இளநீர்} (iḷanīr)}{இளநீர் (iḷanīr)}]{\ta{இளநீர்} (iḷanīr) — tender coconut water}

\label{lesson:TA-C214-ilanir}



\begin{quote}
Before the new one: say the Tamil for buttermilk, then the Tamil for fruit juice.
\end{quote}

\subsection*{You'll want to know: \ta{இளநீர்}}
\textbf{\ta{இளநீர்}} — \emph{iḷanīr} — \textquotedblleft{}tender coconut water\textquotedblright{}.

\begin{grammarlens}[title={What you've built}]
One more word for this chapter.
\end{grammarlens}

\subsection*{Your turn}
\begin{itemize}
\raggedright
  \item \emph{Say it:} \emph{iḷanīr}
  \item \emph{Say it:} \emph{iḷanīr}, once more
  \item \emph{Say it:} say \emph{iḷanīr}
  \item \emph{Recall:} say the Tamil for buttermilk, then the Tamil for fruit juice, then say \emph{iḷanīr} again
\end{itemize}

\begin{tcolorbox}[breakable,colback=teal!4,colframe=teal!35!black,title={Before you move on}]
What does \ta{இளநீர்} mean? (\textquotedblleft{}Tender coconut water\textquotedblright{}.) Say it once more.
\end{tcolorbox}
\section[\texorpdfstring{\ta{மாவு} (māvu)}{மாவு (māvu)}]{\ta{மாவு} (māvu) — flour; batter}

\label{lesson:TA-C214-mavu}



\begin{quote}
Before the new one: say the Tamil for fruit juice, then the Tamil for tender coconut water.
\end{quote}

\subsection*{You'll want to know: \ta{மாவு}}
\textbf{\ta{மாவு}} — \emph{māvu} — \textquotedblleft{}flour; batter\textquotedblright{}.

\begin{grammarlens}[title={What you've built}]
One more word for this chapter.
\end{grammarlens}

\subsection*{Your turn}
\begin{itemize}
\raggedright
  \item \emph{Say it:} \emph{māvu}
  \item \emph{Say it:} \emph{māvu}, once more
  \item \emph{Say it:} say \emph{māvu}
  \item \emph{Recall:} say the Tamil for fruit juice, then the Tamil for tender coconut water, then say \emph{māvu} again
\end{itemize}

\begin{tcolorbox}[breakable,colback=teal!4,colframe=teal!35!black,title={Before you move on}]
What does \ta{மாவு} mean? (\textquotedblleft{}Flour; batter\textquotedblright{}.) Say it once more.
\end{tcolorbox}
\section[\texorpdfstring{\ta{ரொட்டி} (roṭṭi)}{ரொட்டி (roṭṭi)}]{\ta{ரொட்டி} (roṭṭi) — bread}

\label{lesson:TA-C214-rotti}



\begin{quote}
Before the new one: say the Tamil for tender coconut water, then the Tamil for flour; batter.
\end{quote}

\subsection*{You'll want to know: \ta{ரொட்டி}}
\textbf{\ta{ரொட்டி}} — \emph{roṭṭi} — \textquotedblleft{}bread\textquotedblright{}.

\begin{grammarlens}[title={What you've built}]
One more word for this chapter.
\end{grammarlens}

\subsection*{Your turn}
\begin{itemize}
\raggedright
  \item \emph{Say it:} \emph{roṭṭi}
  \item \emph{Say it:} \emph{roṭṭi}, once more
  \item \emph{Say it:} say \emph{roṭṭi}
  \item \emph{Recall:} say the Tamil for tender coconut water, then the Tamil for flour; batter, then say \emph{roṭṭi} again
\end{itemize}

\begin{tcolorbox}[breakable,colback=teal!4,colframe=teal!35!black,title={Before you move on}]
What does \ta{ரொட்டி} mean? (\textquotedblleft{}Bread\textquotedblright{}.) Say it once more.
\end{tcolorbox}
